\documentclass[11pt,a4paper]{article}

\usepackage[utf8]{inputenc}
\usepackage[portuguese]{babel}
\usepackage{url}
\usepackage{xcolor}
\usepackage{listings}

\lstdefinestyle{command}{
    basicstyle=\ttfamily\small,
    backgroundcolor=\color{gray!10},
    frame=single,
    framerule=0.3pt,
    rulecolor=\color{gray!50},
    xleftmargin=0.8cm,
    xrightmargin=0.8cm,
    breaklines=true,
    columns=fullflexible
}

\title{Descrição da Organização}
\author{Grupo 20}
\date{Junho de 2026}

\begin{document}

\maketitle

\section{Identificação do projeto}

O projeto \textit{Insurance Portal Workflow} foi desenvolvido no âmbito da unidade curricular de Projeto e Seminário da Licenciatura em Engenharia Informática e de Computadores, no Instituto Superior de Engenharia de Lisboa, durante o semestre de verão de 2025/2026.

O trabalho foi realizado pelo Grupo 20, constituído pelos seguintes elementos:

\begin{itemize}
  \item 46362 -- Eduardo Franco Bezerra;
  \item 51618 -- Francisco Duarte Tavares;
  \item 51870 -- Rafael Henrique Souza Reis.
\end{itemize}

O projeto foi orientado pelos docentes Artur Jorge Ferreira e Pedro Miguel Florindo Miguens Matutino.

\section{Repositório do projeto}

O código-fonte e a documentação do projeto encontram-se disponíveis no repositório GitHub do projeto:

\begin{center}
  \url{https://github.com/Eduubez/IPW}
\end{center}

O repositório contém a implementação da aplicação, os scripts de base de dados, a configuração dos \textit{containers} Docker e a documentação produzida ao longo do desenvolvimento.

Para a versão beta, o desenvolvimento principal da aplicação encontra-se na \texttt{tag} v1.0.

\section{Organização do repositório}

O repositório encontra-se organizado por módulos e diretorias, separando o código da aplicação, os scripts de base de dados, a configuração dos \textit{containers} e a documentação do projeto.

As principais diretorias são as seguintes:

\begin{itemize}
  \item \texttt{modules-backend} -- contém o código do backend da aplicação, desenvolvido em Kotlin com Spring Boot. Este módulo inclui a camada HTTP, os serviços, o domínio, os repositórios e a configuração de acesso à base de dados;
  \item \texttt{js} -- contém o código do frontend, desenvolvido em React com TypeScript;
  \item \texttt{docker} -- contém ficheiros de apoio à execução da aplicação em \textit{containers}, incluindo os scripts de inicialização da base de dados PostgreSQL;
  \item \texttt{SQL} -- contém os scripts SQL utilizados durante o desenvolvimento, incluindo a criação do esquema, limpeza da base de dados e inserção de dados iniciais;
  \item \texttt{Docs} -- contém a documentação do projeto, incluindo o relatório, imagens e documentos associados às entregas.
\end{itemize}

Na raiz do repositório encontram-se também ficheiros de configuração relevantes, como o \texttt{docker-compose.yml}, utilizado para executar os serviços da aplicação. Existem ainda ficheiros \texttt{Dockerfile} específicos para a construção dos \textit{containers} do backend e do frontend.

\section{Pré-requisitos}

Para executar a versão beta da aplicação é necessário ter instalado:

\begin{itemize}
  \item \textit{Git}, para obter o código-fonte a partir do repositório;
  \item \textit{Docker Desktop}, para construir e executar os \textit{containers} da aplicação;
  \item um navegador web, para aceder à interface da aplicação.
\end{itemize}

Não é necessária a instalação manual de PostgreSQL, Node.js ou Gradle para a execução normal da versão beta, uma vez que estes componentes são utilizados através dos \textit{containers} definidos no ficheiro \texttt{docker-compose.yml}.

\section{Execução da aplicação}

Após obter acesso ao repositório, a aplicação pode ser executada localmente recorrendo ao Docker.

\begin{enumerate}
  \item Fazer o clone do repositório:

        \begin{lstlisting}[style=command]
git clone https://github.com/Eduubez/IPW.git
        \end{lstlisting}

  \item Fazer checkout para a seguinte tag:

        \begin{lstlisting}[style=command]
git checkout v1.0
        \end{lstlisting}

  \item Aceder à diretoria raiz do projeto:

        \begin{lstlisting}[style=command]
cd IPW
        \end{lstlisting}

  \item Construir e iniciar todos os serviços necessários à execução da aplicação:

        \begin{lstlisting}[style=command]
docker compose up -d --build
        \end{lstlisting}
\end{enumerate}

Este comando inicia os serviços definidos no ficheiro \texttt{docker-compose.yml}, incluindo a base de dados PostgreSQL, o backend, o frontend e o serviço MinIO utilizado para armazenamento de anexos.

Após a inicialização dos serviços, os principais pontos de acesso são:

\begin{itemize}
  \item Interface da aplicação: \url{http://localhost:5173};
  \item Backend: \url{http://localhost:8080};
  \item Consola web do MinIO: \url{http://localhost:9001}.
\end{itemize}

Para terminar a execução da aplicação, pode ser usado:

\begin{lstlisting}[style=command]
docker compose down
    \end{lstlisting}

Caso se pretenda remover também os volumes associados aos serviços, incluindo os dados persistidos localmente, pode ser usado:

\begin{lstlisting}[style=command]
docker compose down -v
    \end{lstlisting}

\section{Dados iniciais e utilizadores de teste}

Durante a inicialização da base de dados são executados scripts SQL que criam o esquema da aplicação e inserem dados iniciais para demonstração. Estes dados incluem áreas, utilizadores, papéis e processos de exemplo.

Todos os utilizadores de teste têm a password \texttt{12345}. Alguns utilizadores disponíveis na versão beta são:

\begin{itemize}
  \item \texttt{alice@ipw.pt} -- utilizador com papel de triador;
  \item \texttt{bob@ipw.pt} -- utilizador com papel de averiguador;
  \item \texttt{carol@ipw.pt} -- utilizador com papel de supervisor;
  \item \texttt{david@ipw.pt} -- utilizador com papel de gestor;
  \item \texttt{eve@ipw.pt} -- utilizador com papel de administrador.
\end{itemize}

Estes utilizadores permitem testar os principais fluxos da aplicação, incluindo autenticação, seleção de papel, consulta de processos e criação ou acompanhamento de processos, de acordo com as permissões associadas a cada papel.

\section{Notas sobre a versão beta}

A versão beta disponibilizada tem como objetivo demonstrar os principais fluxos da aplicação, incluindo autenticação, seleção de papel, consulta de processos, criação de processos e acompanhamento do seu estado.

Algumas funcionalidades encontram-se ainda em desenvolvimento ou poderão sofrer alterações até à versão final.

\end{document}
